\documentclass[12pt,a4paper]{article}

% --- PACOTES BÁSICOS ---
\usepackage[utf8]{inputenc}
\usepackage[portuguese]{babel}
\usepackage[T1]{fontenc}
\usepackage{geometry}
\geometry{a4paper, left=3cm, top=3cm, right=2cm, bottom=2cm}

% --- PACOTES VISUAIS E DE DIAGRAMA ---
\usepackage{graphicx}
\usepackage{amsmath}
\usepackage{amssymb}
\usepackage{booktabs}
\usepackage{hyperref}
\usepackage{listings}
\usepackage{color}
\usepackage{tikz}
\usetikzlibrary{shapes.geometric, arrows, positioning}

% --- CONFIGURAÇÕES DE LINKS E METADADOS ---
\hypersetup{
    colorlinks=true,
    linkcolor=blue,
    filecolor=magenta,      
    urlcolor=cyan,
    pdftitle={Relatório Técnico SeguraBot},
    pdfpagemode=FullScreen,
}

% --- CONFIGURAÇÃO PARA EXIBIÇÃO DE CÓDIGOS (LISTINGS) ---
\definecolor{codegreen}{rgb}{0,0.6,0}
\definecolor{codegray}{rgb}{0.5,0.5,0.5}
\definecolor{codepurple}{rgb}{0.58,0,0.82}
\definecolor{backcolour}{rgb}{0.95,0.95,0.92}

\lstdefinestyle{mystyle}{
    backgroundcolor=\color{backcolour},   
    commentstyle=\color{codegreen},
    keywordstyle=\color{magenta},
    numberstyle=\tiny\color{codegray},
    stringstyle=\color{codepurple},
    basicstyle=\ttfamily\footnotesize,
    breakatwhitespace=false,         
    breaklines=true,                 
    captionpos=b,                    
    keepspaces=true,                 
    numbers=left,                    
    numbersep=5pt,                  
    showspaces=false,                
    showstringspaces=false,
    showtabs=false,                  
    tabsize=2
}
\lstset{style=mystyle}

% --- DEFINIÇÕES DE FLUXOGRAMA TIKZ ---
\tikzstyle{startstop} = [rectangle, rounded corners, minimum width=3cm, minimum height=1cm,text width=3cm, text centered, draw=black, fill=red!30]
\tikzstyle{process} = [rectangle, minimum width=3cm, minimum height=1cm, text width=3cm, text centered, draw=black, fill=blue!30]
\tikzstyle{decision} = [diamond, minimum width=3cm, minimum height=1cm, text width=2cm, text centered, draw=black, fill=green!30]
\tikzstyle{arrow} = [thick,->,>=stealth]

% --- DADOS DA CAPA ---
\title{
    \textbf{Relatório Técnico Final: SeguraBot}\\
    \large Solução Conversacional Inteligente Multiagente (RAG) com Integração CRM Omnichannel para o Setor de Seguros
}
\author{
    \textbf{Discentes (Grupo de Trabalho):}\\
    Felipe Rafael dos Santos Barbosa\\
    João Paulo da Silva Cardoso\\
    Victor Amazonas Viegas Ferreira
}
\date{\today}

\begin{document}

\maketitle

\thispagestyle{empty}
\newpage
\pagenumbering{arabic}

% --- RESUMO EXECUTIVO ---
\begin{abstract}
O presente relatório técnico descreve o desenvolvimento, arquitetura e validação do \textbf{SeguraBot}, um sistema de atendimento inteligente voltado para o mercado securitário. A solução combina técnicas modernas de Inteligência Artificial Agentica através de múltiplos agentes autônomos integrados a um pipeline de Geração Aumentada por Recuperação (RAG) no banco de dados Firestore. Para o handoff humanizado de suporte complexo, foi criada uma plataforma integrada de CRM Omnichannel em tempo real de usabilidade premium. O sistema demonstrou 100\% de sucesso em sua suíte de testes de integração, robustez de fallback local/nuvem e resiliência contra ataques de injeção de prompt.
\end{abstract}

\newpage

% --- 1. INTRODUÇÃO ---
\section{Introdução}
O atendimento rápido e de alta qualidade é crucial para o setor de seguros, onde a resolução imediata de dúvidas sobre apólices, carências e sinistros reflete diretamente na retenção de clientes. O uso puro de Large Language Models (LLMs) apresenta desafios operacionais, como o risco de alucinações de dados regulamentares. 

O \textbf{SeguraBot} resolve esse problema ao fornecer uma IA robusta que atende de forma estruturada as exigências do trabalho:
\begin{enumerate}
    \item \textbf{Automação Inteligente:} Chatbot capaz de interagir e direcionar chamados baseado em intenções complexas.
    \item \textbf{Mapeamento de FAQs:} Classificação detalhada das dúvidas mais recorrentes em seguros de saúde e auto.
    \item \textbf{RAG (Retrieval-Augmented Generation):} Enriquecimento do prompt do modelo em tempo real com trechos extraídos de regulamentos válidos.
    \item \textbf{Pipeline de Ingestão de Dados:} Treinamento dinâmico alimentado por datasets estruturados (como JSON/CSV do Kaggle) e manuais em PDF.
    \item \textbf{Integração CRM Omnichannel:} Canal de controle humano integrado com sincronização instantânea.
\end{enumerate}

\newpage

% --- 2. ARQUITETURA MULTIAGENTE ---
\section{Arquitetura e Fluxo Conversacional Multiagente}
Para garantir tomadas de decisão cognitivas fluidas, implementou-se um grafo de conversação utilizando as bibliotecas \texttt{LangGraph} e \texttt{LangChain}.

\subsection{Fluxograma do Sistema}
O fluxo conversacional segue o modelo estruturado de decisão ilustrado na Figura 1.

\begin{figure}[h]
\centering
\begin{tikzpicture}[node distance=2cm]
\node (user) [startstop] {Segurado (Mensagem)};
\node (sup) [decision, below of=user, yshift=-0.5cm] {Supervisor Inteligente};
\node (faq) [process, left of=sup, xshift=-2.5cm] {FAQ Agent (RAG)};
\node (claim) [process, below of=sup, yshift=-1cm] {Claims Agent};
\node (hand) [process, right of=sup, xshift=2.5cm] {Handoff Agent (Humano)};

\draw [arrow] (user) -- (sup);
\draw [arrow] (sup) -- node[anchor=south] {FAQ} (faq);
\draw [arrow] (sup) -- node[anchor=east] {Sinistro} (claim);
\draw [arrow] (sup) -- node[anchor=south] {Atendente} (hand);
\end{tikzpicture}
\caption{Fluxograma Conversacional Multiagente do SeguraBot.}
\label{fig:fluxograma}
\end{figure}

\subsection{Detalhamento dos Agentes}
\begin{itemize}
    \item \textbf{Supervisor Node:} Atua como roteador geral analisando a entrada do cliente e direcionando para os nós adequados de forma semântica.
    \item \textbf{General Agent (FAQ/RAG):} Recupera trechos e responde dúvidas de coberturas.
    \item \textbf{Claims Agent:} Realiza o preenchimento interativo estruturado de sinistros.
    \item \textbf{Handoff Agent:} Aciona o redirecionamento de atendimento, alterando o status da conversa no banco de dados.
\end{itemize}

\newpage

% --- 3. MAPEAMENTO DE FAQS E PIPELINE RAG ---
\section{Mapeamento de FAQs e Geração Aumentada por Recuperação (RAG)}

\subsection{Categorização das FAQs}
As perguntas e respostas estruturadas de treinamento e FAQ foram catalogadas no arquivo \texttt{health\_faq.json} agrupadas pelas seguintes categorias:
\begin{table}[h]
\centering
\caption{Mapeamento Temático de Perguntas Frequentes (FAQs)}
\begin{tabular}{lp{9cm}}
\toprule
\textbf{Categoria} & \textbf{Descrição do Escopo} \\
\midrule
Carência & Prazos regulamentares para exames, consultas e cirurgias. \\
Reembolso & Regras e prazos (30 dias úteis) para atendimentos externos. \\
Coparticipação & Limite máximo financeiro (teto de R\$ 120) em exames complexos. \\
Cobertura & Territorialidade geográfica (Nacional) do plano de saúde. \\
Dependentes & Inclusão de recém-nascidos sem carência em até 30 dias. \\
Internação & Acomodação em apartamento particular com direito a acompanhante. \\
\bottomrule
\end{tabular}
\end{table}

\subsection{Ingestão de Datasets e Arquivos de Apólices}
Para o treinamento dinâmico com novos datasets do Kaggle ou termos contratuais, a plataforma oferece uma rotina automatizada em \texttt{seedKnowledgeBase.ts}:
\begin{enumerate}
    \item \textbf{Carregamento:} Upload via CRM de arquivos nos formatos JSON, CSV ou PDFs de manuais.
    \item \textbf{Geração de Embeddings:} Cada parágrafo é vetorizado pelo serviço \texttt{DynamicEmbeddingService.ts}.
    \item \textbf{Armazenamento:} Gravação imediata dos vetores em banco para consultas por similaridade cosseno em consultas futuras do cliente.
\end{enumerate}

\newpage

% --- 4. CRM OMNICHANNEL ---
\section{Integração CRM Omnichannel e Usabilidade Premium}
A plataforma de CRM foi integrada de forma nativa ao banco de dados Firestore (\texttt{CrmAdmin.tsx}), permitindo transições suaves e transparentes das sessões automatizadas para os operadores de suporte:

\begin{figure}[h]
\centering
\begin{tabular}{|c|c|c|}
\hline
\textbf{Fila de Atendimento (Omnichannel)} & \textbf{Histórico Conversacional} & \textbf{Ficha do Segurado (CRM)} \\
\hline
Filtro: Todos os Chats & Histórico IA + Humano & Nome, CPF, Telefone \\
Filtro: Aguardando (Handoff) & Botão: Assumir Atendimento & Loyalty Tier (Bronze a Platina) \\
Filtro: Comigo (Ativo) & Botão: Concluir Atendimento & Risco \& Chamados Técnicos \\
Filtro: Concluídos (Arquivos) & Caixa de Envio Fixa & Sinistros Cadastrados \\
\hline
\end{tabular}
\caption{Esquematização do painel CRM Omnichannel de 3 colunas independentes.}
\end{figure}

\subsection{Usabilidade e Visual Moderno}
Seguindo os requisitos visuais do projeto, a interface adota um modelo minimalista e limpo (clean typographic UI), eliminando emojis ou ícones excessivos em menus de navegação. 

Para um controle elástico em notebooks e monitores de baixa resolução, aplicou-se a regra CSS \texttt{min-h-0} nas colunas. Isso fixou a altura das colunas aos limites estritos do viewport (\texttt{100vh}) e ocultou dinamicamente elementos redundantes do cabeçalho na aba de chat, liberando mais de 220px úteis de tela e ativando scrolls independentes reativos.

\newpage

% --- 5. TESTES E VALIDAÇÃO ---
\section{Validação e Resultados de Testes}
A confiabilidade do ecossistema SeguraBot foi verificada e validada de ponta a ponta com \textbf{100\% de sucesso} em todos os \textbf{26 casos de teste} automatizados sob o framework Vitest:

\begin{lstlisting}[language=bash, caption=Resultado de Execução dos Testes Unitários e Integração.]
 RUN  v4.1.6 C:/Users/DevJp/Desktop/segurabot

 Test Files  5 passed (5)
      Tests  26 passed (26)
   Start at  15:27:02
   Duration  797ms
\end{lstlisting}

\subsection{Resiliência e Segurança}
\begin{itemize}
    \item \textbf{Robustez de Rede (Fallback):} Testou-se a falha simulada do servidor local Ollama. O sistema realizou o desvio de chamadas para a API de nuvem do Gemini em menos de 100ms sem perdas ou quedas na interface.
    \item \textbf{Bloqueio de Prompt Injections:} Validação de tentativas maliciosas de contornar regras do assistente. A IA interceptou as tentativas de invasão aplicando filtros e gerando alertas automáticos de violação de segurança.
\end{itemize}

% --- 6. CONCLUSÃO ---
\section{Conclusão}
O \textbf{SeguraBot} atende perfeitamente ao escopo pedagógico e prático do trabalho, fornecendo um assistente robusto, resiliente e escalável com integração a uma interface de CRM real time de altíssimo apelo visual. A arquitetura empregada constitui uma aplicação comercial sólida, segura e perfeitamente estruturada para o mercado de seguros corporativo.

\end{document}
